\chapter{Cancer Immunotherapy}
\index{Cancer Immunotherapy}

\section*{Definition and Overview}
Cancer immunotherapy is a type of cancer treatment that helps the body's own immune system recognise and attack cancer cells. Unlike chemotherapy, which acts directly on rapidly dividing cells, immunotherapy works by boosting, directing, or unleashing the immune response so that the immune system itself destroys the tumour.

Treatment includes immune checkpoint inhibitors, monoclonal antibodies, cancer vaccines, and cell-based treatments such as CAR-T cell therapy. Immunotherapy has changed the outlook for some cancers that were previously very difficult to treat, including advanced melanoma and certain lung, kidney, and bladder cancers, and it is now a standard part of care for many tumour types.

\section*{Epidemiology}
Immunotherapy is now a routine part of treatment for many cancers and is used increasingly around the world. Melanoma, lung cancer, kidney cancer, bladder cancer, and Hodgkin lymphoma are among the cancers most responsive to checkpoint inhibitors, and these drugs are now used in the first-line treatment of some advanced cancers. Because many cancers are diagnosed at a late stage, immunotherapy has grown rapidly over the past decade and is one of the fastest-developing areas of cancer care. Its expanding use reflects both proven benefit and continuing research into new agents and combinations.

\section*{Aetiology and Pathophysiology}
Cancer cells arise from the body's own tissues and can evade the immune system, which normally destroys abnormal cells. Immunotherapy targets the specific ways tumours escape immune attack:
\begin{itemize}
    \item \textbf{Immune evasion:} many cancers display checkpoints, such as PD-L1, that switch off T-cells and prevent an immune attack
    \item \textbf{Checkpoint inhibitors:} drugs that block these checkpoints (PD-1, PD-L1, and CTLA-4), freeing T-cells to attack the tumour
    \item \textbf{Monoclonal antibodies:} engineered antibodies that mark cancer cells for destruction or block growth signals
    \item \textbf{Cell therapy (CAR-T):} a patient's own immune cells are collected, genetically engineered to recognise the cancer, and returned to destroy it
    \item \textbf{Cancer vaccines:} designed to train the immune system against tumour-specific proteins
    \item \textbf{Tumour dependence:} cancers with many mutations or high PD-L1 expression often respond best
\end{itemize}

\section*{Clinical Features}
Immunotherapy is delivered as a treatment, and its clinical features are those of the response and the side effects:
\begin{itemize}
    \item \textbf{Response:} tumours may shrink or stabilise, sometimes over months, and responses can be long-lasting
    \item \textbf{Treatment delivery:} usually given as an intravenous infusion every few weeks in a hospital or clinic
    \item \textbf{Flu-like symptoms:} fever, chills, and fatigue are common around the time of infusion
    \item \textbf{Autoimmune-like side effects:} because immunotherapy activates the immune system, it can cause inflammation of normal organs, including the skin, bowel, liver, lungs, and thyroid
    \item \textbf{Treatment holidays:} some people are able to stop treatment after a period of disease control
\end{itemize}

\section*{Differential Diagnosis}
Side effects of immunotherapy can resemble other problems and need careful evaluation:
\begin{itemize}
    \item Diarrhoea from immune-related colitis versus infection or the cancer itself
    \item A new cough or breathlessness from immune-related lung inflammation (pneumonitis) versus infection, a blood clot, or worsening cancer
    \item Fatigue from treatment versus anaemia, underactive thyroid, or depression
    \item A rash from the drug versus an allergic reaction or another skin condition
    \item Headache from immune-related brain inflammation versus another cause
\end{itemize}

\section*{When to Seek Medical Care}
Any new or worsening symptom while on immunotherapy should be reported to the cancer care team promptly, because immune-related side effects are easier to treat when caught early.

\textbf{Contact your local emergency services for:}
\begin{itemize}
    \item Difficulty breathing, chest pain, or coughing up blood
    \item Severe abdominal pain or diarrhoea with blood
    \item Sudden weakness, confusion, or seizures
    \item A high fever or feeling very unwell, which may indicate a serious infection
    \item Symptoms of a severe allergic reaction, such as swelling of the face or throat
\end{itemize}

\section*{Diagnosis and Investigations}
Immunotherapy is selected on the basis of the cancer type, and side effects are diagnosed and monitored with targeted tests:
\begin{itemize}
    \item \textbf{Tumour testing:} the tumour may be tested for markers such as PD-L1, and genetic testing may assess the tumour mutational burden
    \item \textbf{Imaging:} CT, PET, or MRI scans are used to assess the tumour response
    \item \textbf{Blood tests:} regular full blood counts and liver, kidney, and thyroid tests monitor for immune-related effects
    \item \textbf{Response assessment:} some apparent growth on scans is caused by inflammation and can be misleading, so responses are judged over time
    \item \textbf{Special tests:} if side effects are suspected, tests such as lung imaging, stool tests, or skin biopsy may be used
\end{itemize}

\section*{Management}
Management is led by a cancer team and combines the treatment itself with careful monitoring and support:
\begin{itemize}
    \item \textbf{Treatment plan:} immunotherapy may be given alone or with chemotherapy, targeted therapy, or radiation
    \item \textbf{Immune-related side effects:} mild effects may settle with observation or simple treatment, while more serious ones are treated with corticosteroids or stronger immune-suppressing drugs, often with a pause in immunotherapy
    \item \textbf{Supportive care:} management of fatigue, nutrition, and pain, together with psychological support throughout treatment
    \item \textbf{Regular monitoring:} frequent blood tests, scans, and clinic reviews to track both the cancer and the side effects
    \item \textbf{Stopping treatment:} immunotherapy may be stopped after a good response or if side effects become severe
    \item \textbf{Clinical trials:} access to newer immunotherapy approaches may be offered through research studies
\end{itemize}

\section*{Complications}
\begin{itemize}
    \item Immune-related inflammation of any organ, most commonly the skin, bowel, liver, thyroid, and lungs, and less often the heart, nerves, or kidneys
    \item Serious autoimmune reactions that can be life-threatening if not recognised and treated early
    \item Infusion reactions, including allergic reactions during or shortly after treatment
    \item Infection, especially in people also receiving chemotherapy that lowers white cell counts
    \item The treatment may not work, or the cancer may eventually become resistant
    \item Long-term effects of immune activation, including underactive thyroid and other gland problems
\end{itemize}

\section*{Prognosis and Outcomes}
Immunotherapy can produce durable responses, and in some advanced cancers -- especially melanoma and certain lung cancers -- a significant minority of patients survive for many years, which was rare before these treatments existed. However, responses vary widely: some tumours shrink dramatically, others are unaffected, and some respond briefly before growing again. Side effects are common but, when caught early, are usually manageable. The outlook depends on the cancer type, its stage, and how well the person tolerates treatment.

\section*{Prevention and Self-Care}
\begin{itemize}
    \item Attend every infusion and blood test as scheduled, so that side effects are caught early
    \item Report any new symptom, however mild, to your cancer team straight away, including rash, diarrhoea, cough, headache, or feeling unwell
    \item Keep a symptom diary and take it to appointments
    \item Eat well, stay as active as you can, and rest when you need to
    \item Look after your skin and protect it from the sun
    \item Let every doctor and pharmacist you see know that you are receiving immunotherapy
    \item Seek support from family, friends, and cancer support services for the emotional side of treatment
\end{itemize}

\section*{Sources and Further Reading}
The following sources provide reliable, accessible information on cancer immunotherapy:
\begin{itemize}
    \item National Cancer Institute. \textit{Immunotherapy to treat cancer.} https://www.cancer.gov/about-cancer/treatment/
    \item Cancer Council. \textit{Immunotherapy.} https://www.cancer.org.au/
    \item NHS. \textit{Immunotherapy.} https://www.nhs.uk/conditions/
    \item American Cancer Society. \textit{Immunotherapy.} https://www.cancer.org/
    \item Mayo Clinic. \textit{Immunotherapy for cancer.} https://www.mayoclinic.org/
    \item MSD Manual. \textit{Biology of the immune system in cancer.} https://www.msdmanuals.com/
\end{itemize}
